% !TeX program = xelatex
% !TeX options = -shell-escape
% !TEX encoding = UTF-8

% --------------------------------------------------------------
% This is all preamble stuff that you don't have to worry about.
% Head down to where it says "Start here"
% --------------------------------------------------------------

\documentclass[12pt]{article}
\usepackage[margin=2cm]{geometry} % 頁面邊距設置
\usepackage{amsmath,amssymb,amsthm} % 數學相關套件
\usepackage{graphicx} % 圖形插入
\usepackage{booktabs} % 表格線條
\usepackage[AutoFakeBold,AutoFakeSlant]{xeCJK} % 中文支援
\usepackage{fancyhdr} % 頁首頁尾設置
\usepackage{xcolor,changepage,mdframed} % 顏色與框線相關
\usepackage{titlesec} % 標題格式
\usepackage{setspace} % 行距
\usepackage{minted} % 程式碼高亮（需安裝 Python 套件 pygments）
\usepackage{booktabs}
\usepackage{caption}
\usepackage{array}
\usepackage{ulem}
\usepackage[shortlabels]{enumitem} % 列表格式
\usepackage[colorlinks=true,linkcolor=black,urlcolor=blue,citecolor=blue]{hyperref}
% 超連結
\graphicspath{{images}}

% 字體設定 (xeCJK)
\setCJKmainfont{黑體-繁}
\setCJKmonofont{Noto Sans Mono CJK TC}

% 頁首/頁尾設定 (fancyhdr)
\pagestyle{fancy}
\fancyhead[L]{NASA Homework 12} % 可以用 \leftmark 取代
\fancyhead[R]{B13901011 吳亞倫}
\fancyfoot[C]{\thepage}
\renewcommand{\headrulewidth}{0.4pt}
\renewcommand{\footrulewidth}{0pt}
\setlength{\headheight}{15pt}

% tex-fmt: off
\newtheorem*{theorem}{Theorem}
\newenvironment{lemma}[2][Lemma]{\begin{trivlist}
\item[\hskip \labelsep {\bfseries #1}\hskip \labelsep {\bfseries #2.}]}{\end{trivlist}}
\newenvironment{exercise}[2][Exercise]{\begin{trivlist}
\item[\hskip \labelsep {\bfseries #1}\hskip \labelsep {\bfseries #2.}]}{\end{trivlist}}
\newenvironment{problem}[2][Problem]{\begin{trivlist}
\item[\hskip \labelsep {\bfseries #1}\hskip \labelsep {\bfseries #2}]}{\end{trivlist}}
\newenvironment{question}[2][Question]{\begin{trivlist}
\item[\hskip \labelsep {\bfseries #1}\hskip \labelsep {\bfseries #2.}]}{\end{trivlist}}
\newenvironment{corollary}[2][Corollary]{\begin{trivlist}
\item[\hskip \labelsep {\bfseries #1}\hskip \labelsep {\bfseries #2.}]}{\end{trivlist}}
\newenvironment{solution}{\begin{proof}[Solution]}{\end{proof}}
\newenvironment{answer}{\begin{proof}[Ans]}{\end{proof}}
\renewcommand{\thesubsection}{\thesection. (\alph{subsection})}

% 樣式設定

\setminted{frame=lines, linenos, framesep=2mm, breaklines=true, breakanywhere=true, autogobble=true}

% 自訂框線環境
\newmdenv[
  linewidth=2pt, linecolor=gray, backgroundcolor=gray!10,
  topline=false, bottomline=false, rightline=false,
  skipabove=10pt, skipbelow=10pt, innerleftmargin=10pt,
  innerrightmargin=0pt, innertopmargin=5pt, innerbottommargin=5pt
]{myquote}

\newenvironment{mycolorbox}[1][gray]{
  \renewmdenv[
    linewidth=2pt, linecolor=#1, backgroundcolor=#1!10,
    topline=false, bottomline=false, rightline=false,
    skipabove=10pt, skipbelow=10pt, innerleftmargin=10pt,
    innerrightmargin=0pt, innertopmargin=5pt, innerbottommargin=5pt
  ]{myquote}
  \begin{myquote}
}{\end{myquote}}

\linespread{1.1}
% \usemintedstyle{xcode}
\AtBeginEnvironment{minted}{\vspace{-10pt}} % 可根據需求調整間距數值
\setminted{breaklines=true, breakanywhere=true, autogobble=true}

% tex-fmt: on
\begin{document}
\thispagestyle{fancy}

%------------------------------------------------------------
%                         Start here
% -----------------------------------------------------------

\begin{center}
  {\large \bfseries Network Administration / System Administration} \\
  \vspace{1em}
  {\Large \bfseries Homework 12} \\
\end{center}

\vspace{2em}

\begin{tabular}{@{}l l}
  \textbf{Student:} & 吳亞倫 \\
  \textbf{ID:} & B13901011  \\
  \textbf{Department:} & 電機工程學系
\end{tabular}
\bigskip

\begin{mycolorbox}[cyan]
  \textbf{Discussion Partners:}
  \begin{itemize}
    \item Classmates：賴泓安、鄭竣陽、喻慈恩
      \vspace{0.5em}
  \end{itemize}
\end{mycolorbox}

\section{Linux 大小事}
\begin{mycolorbox}[teal]
  \textbf{Reference:}
  \footnotesize
  \begin{itemize}
    \item 
  \end{itemize}
\end{mycolorbox}

% Problem 1(a)
\subsection{}

% Problem 1(b)
\subsection{}

% Problem 1(c)
\subsection{}

% Problem 1(d)
\subsection{}

\clearpage
\section{畫中有話}
\begin{mycolorbox}[teal]
  \textbf{Reference:}
  \footnotesize
  \begin{itemize}
    \item \url{https://www.cnblogs.com/znh8/p/9316923.html}
  \end{itemize}
\end{mycolorbox}

\subsection{hide.py 運作原理}

\subsection{圖中的訊息}


\clearpage
\section{Alya Judge}
\begin{mycolorbox}[teal]
  \textbf{Reference:}
  \footnotesize
  \begin{itemize}
    \item 
  \end{itemize}
\end{mycolorbox}

\subsection{fysty 的上傳紀錄}
\textit{Steps:}
\begin{enumerate}
  \item jef;
  \begin{minted}{text}
    http://140.112.91.4:45510/submissions/%2e%2e%2faccounts%2faccounts
  \end{minted}
\end{enumerate}

\subsection{未完成的題目公開了}

\subsection{admin 的上傳紀錄}



% -----------------------------------------------------------
%    You don't have to mess with anything below this line.
% -----------------------------------------------------------

\end{document}
